\documentclass[11en, a4paper, oneside]{article}
\usepackage[utf8]{inputenc}   % <<<<< Linux
\usepackage[english]{babel} % <<<<< English
\usepackage{notoccite}
\usepackage[skip=0.5\baselineskip]{caption}
\hyphenation{GTKWave}
\usepackage{listings}
\usepackage[all]{nowidow}
\usepackage{tabularx}
\usepackage{amsmath}
\usepackage{float}

%blind text
\usepackage{lipsum}

\usepackage{graphicx}
\graphicspath{ {./} {../../figlib/} }
\def\FontLn{% 16 pt normal
  \usefont{T1}{phv}{m}{n}\fontsize{16pt}{16pt}\selectfont}
\def\FontLb{% 16 pt bold
  \usefont{T1}{phv}{b}{n}\fontsize{16pt}{16pt}\selectfont}
\def\FontMn{% 14 pt normal
  \usefont{T1}{phv}{m}{n}\fontsize{14pt}{14pt}\selectfont}
\def\FontMb{% 14 pt bold
  \usefont{T1}{phv}{b}{n}\fontsize{14pt}{14pt}\selectfont}
\def\FontSn{% 12 pt normal
  \usefont{T1}{phv}{m}{n}\fontsize{12pt}{12pt}\selectfont}

% Use Arial font as default
%
\renewcommand{\rmdefault}{phv}
\renewcommand{\sfdefault}{phv}
\usepackage{geometry}	
\geometry{verbose,tmargin=2.5cm,bmargin=2.5cm,lmargin=2.5cm,rmargin=2.5cm}

%\usepackage{setspace}
%\renewcommand{\baselinestretch}{1.5}

\usepackage[pdftex]{hyperref} % enhance documents that are to be
                              % output as HTML and PDF
\hypersetup{colorlinks,       % color text of links and anchors,
                              % eliminates borders around links
%            linkcolor=red,    % color for normal internal links
            linkcolor=black,  % color for normal internal links
            anchorcolor=black,% color for anchor text
%            citecolor=green,  % color for bibliographical citations
            citecolor=black,  % color for bibliographical citations
%            filecolor=magenta,% color for URLs which open local files
            filecolor=black,  % color for URLs which open local files
%            menucolor=red,    % color for Acrobat menu items
            menucolor=black,  % color for Acrobat menu items
%            pagecolor=red,    % color for links to other pages
            pagecolor=black,  % color for links to other pages
%            urlcolor=cyan,    % color for linked URLs
            urlcolor=black,   % color for linked URLs
	          bookmarks=true,         % create PDF bookmarks
	          bookmarksopen=false,    % don't expand bookmarks
	          bookmarksnumbered=true, % number bookmarks
	          pdftitle={report},
            pdfauthor={Andre C. Marta},
%            pdfsubject={Thesis Title},
%            pdfkeywords={Thesis Keywords},
            pdfstartview=FitV,
            pdfdisplaydoctitle=true}

\usepackage[numbers,sort&compress]{natbib} % <<<<< References in numbered list [1],[2],...
\usepackage{subcaption} 
\usepackage{mdframed}
\usepackage[utf8]{inputenc}
\usepackage{graphicx}
\usepackage{amsmath}
\usepackage{indentfirst}
\usepackage{fancyhdr}
\usepackage{hyperref}
\usepackage{subfiles}
\usepackage{wrapfig}
\usepackage{subcaption}
\usepackage[utf8]{inputenc}
%\usepackage[nottoc]{tocbibind}
\usepackage{lastpage}
\newcommand{\clap}{\makebox[0pt]}
\usepackage{titlesec}
\usepackage{amsfonts}
\usepackage{listings}
\usepackage{color} %red, green, blue, yellow, cyan, magenta, black, white
\definecolor{mygreen}{RGB}{28,172,0} % color values Red, Green, Blue
\definecolor{mylilas}{RGB}{170,55,241}
\newcommand{\overbar}[1]{\mkern 1.5mu\overline{\mkern-1.5mu#1\mkern-1.5mu}\mkern 1.5mu}
\titleformat{\paragraph}
{\normalfont\normalsize\bfseries}{\theparagraph}{1em}{}
\titlespacing*{\paragraph}
{0pt}{3.25ex plus 1ex minus .2ex}{1.5ex plus .2ex}
\geometry{margin=0.67in}

\usepackage{etoolbox}
\patchcmd{\thebibliography}{\section*{\refname}}{}{}{}

\graphicspath{ {./Imagens/} }

\hypersetup{colorlinks,citecolor=black,filecolor=black,linkcolor=black,urlcolor=black} 

\lstset{language=Matlab,%
    %basicstyle=\color{red},
    breaklines=true,%
    morekeywords={matlab2tikz},
    keywordstyle=\color{blue},%
    morekeywords=[2]{1}, keywordstyle=[2]{\color{black}},
    identifierstyle=\color{black},%
    stringstyle=\color{mylilas},
    commentstyle=\color{mygreen},%
    showstringspaces=false,%without this there will be a symbol in the places where there is a space
    numbers=left,%
    numberstyle={\tiny \color{black}},% size of the numbers
    numbersep=9pt, % this defines how far the numbers are from the text
    emph=[1]{for,end,break},emphstyle=[1]\color{red}, %some words to emphasise
    %emph=[2]{word1,word2}, emphstyle=[2]{style},    
}


\pagestyle{fancy}
\fancyhf{}
\rhead{Circuit Theory and Electronics Fundamentals}
\lhead{Group 21}
\cfoot{Página  \thepage \hspace{1pt} de \pageref{LastPage}}
\pagenumbering{roman}

\begin{document}


\begin{titlepage}
	\begin{center}
		\begin{figure}[htb!]
			\begin{center}
				\includegraphics[width=10cm]{../../figlib/tecnico.png}
			\end{center}
		\end{figure}
		
        \vspace{30pt}
        \begin{center}
        \Large{\center Integrated Masters in Aerospace Engineering, Técnico, University of Lisbon}\\
        \Large{\center Circuit Theory and Electronics Fundamentals}\\
        \end{center}
            
        \vspace{60pt}
        \Huge{\textbf{Laboratory Report 4}}
        
        \vspace{120pt}
        \begin{minipage}{0.4\textwidth}
		\begin{flushleft} \large
			\emph{\LARGE{\textbf{Grupo 21:}}}\par \vspace{10pt}
			95791, Francisco Carvalho \\ \vspace{20pt}
            95805, João Matias\\ \vspace{20pt}
            95846, Simão Gonçalves\\ \vspace{20pt}
		\end{flushleft}
	\end{minipage}
	~
	\begin{minipage}[b]{0.4\textwidth}
		\begin{flushright} \large
        	{}
		\end{flushright}
	\end{minipage}\\[2cm]
       \vspace{10pt}
        \large{April 2021}\\
	\end{center}
\end{titlepage}

\newpage
\renewcommand{\contentsname}{Índice}
\tableofcontents
\thispagestyle{empty}

\newpage
\pagenumbering{arabic}
\setcounter{page}{3}

\section{Introduction}

For the fifth laboratory assignment on our Circuit Theory and Electronics Fundamentals course, we had to dimension and implement a Band Pass Filter (BPF) using an OP-AMP (Operational Amplifier) with a central frequency of 1KHz and a gain at central frequency of 40dB. An OP-AMP  is a DC-coupled high-gain electronic voltage amplifier with a differential input and, usually, a single-ended output. We had available a certain number of components to build the circuit that is shown in the image below. In order to study the behaviour of the circuit, we use Ngspice to measure the output voltage gain in the pass-band, the central frequency, and finally the input and output impedance at the central frequency. Then, we use Octave to do the theoretical analyses (Compute the frequency responses Vo(f)/Vi(f), the gain, and the input and output impedance at the central frequency. In this report, we will perform a comparison between our theoretical analysis and the simulation results, trying to explain any major discrepancies.
\par

To measure the quality of the filter the following equation was used:

\begin{equation}
    Merit=\frac{1}{cost*gaindeviation*centralfreqdeviation + 10^{-6}}
\end{equation}

\par

\begin{figure}[ht] \centering
\includegraphics[width=0.7\linewidth]{circ5.pdf}
\caption{Circuit}
\label{fig:circ}
\end{figure}
\par
A table with the values associated to each component is shown below. Note that the variables are in Volts, Ohms or Faradays.
\par
\begin{table}[H]
\centering
\begin{tabularx}{0.6\textwidth} {
  | >{\raggedright\arraybackslash}X
  | >{\raggedleft\arraybackslash}X | }
 \hline
R1 & 1.000000e+03 \\ \hline
R2 & 1.000000e+03 \\ \hline
R3 & 1.500000e+05 \\ \hline
R4 & 1.000000e+03 \\ \hline
C1 & 2.200000e-07 \\ \hline
C2 & 1.100000e-07 \\ \hline

\end{tabularx}
\caption{Circuit's components}
\end{table}
\clearpage

\section{Theoretical Analysis}

\subsection{Description and Mathematical Considerations}

The Op-Amp (signal amplifier) used was considered ideal, which means that the internal impedance between v+ and v- is infinite and that there is no current flowing through it (v+ = v-). So, by connecting a capacitor ($C_1$) in series with the input voltage, we made a band pass filter, functioning as a high pass filter. Furthermore, in the final stage of the circuit, another capacitor ($C_2$) was connected in parallel with the output voltage, functioning as a low pass filter. Concluding, this circuit consists of a high pass filter, a signal amplifier and a low pass filter in series.

The circuit shown above was analyzed using Octave.

Initially we computed the input and the output impedance:

\begin{equation}
    Z_{in} = R_1 + \frac{1}{j*w*C}
\end{equation}

\begin{equation}
    Z_{out} = \frac{R_2}{j*w_0*C_2*(R_2+ \frac{1}{j w C_2})}
\end{equation}

\begin{table}[H]
\centering
\begin{tabularx}{0.6\textwidth} {
  | >{\raggedright\arraybackslash}X
  | >{\raggedleft\arraybackslash}X | }
 \hline
Input Impedance & 1.000000e+03 + -7.071068e+02j \\ \hline
Output Impedance & 6.666667e+02 + -4.714045e+02j\\ \hline

\end{tabularx}
\caption{Impedances}
\end{table}

Then we computed the gain:

\begin{equation}
    Gain = A_V * A_H * A_L
\end{equation}

\begin{equation}
A_{V}=  1+\frac{R3}{R4}
\end{equation}

\begin{equation}
A_{L}=  \frac{1}{1+R2*C2*s}
\end{equation}

\begin{equation}
A_{H}=  \frac{R1*C1*s}{1+R1*C1*s}
\end{equation}

Where $A_V$ is the gain that results from the OP-AMP (signal amplifier), $A_H$ and $A_L$ are the gains that correspond to the high pass and low pass filters, respectively.

Using the previous equations we obtained the gained presented:

\begin{table}[H]
\centering
\begin{tabularx}{0.6\textwidth} {
  | >{\raggedright\arraybackslash}X
  | >{\raggedleft\arraybackslash}X | }
 \hline
Gain & 1.006667e+02 \\ \hline
Gain & 4.005771e+01 dB \\ \hline

\end{tabularx}
\caption{Gain}
\end{table}

To better understand the frequency response of gain, we plot the following graph (REFERENCIA).  We can see that low and high frequencies have a a low gain and frequencies near to 1000 Hz have the maximum gain, as expected for a band pass filter.
In the first stage the high pass filter blocks the low frequencies and in the final stage we have a low pass filter that blocks high frequencies. To ensure an high gain, we must focus obtaining an high AV , because near the central frequency both AH and AL will be approximately 1. We can see as well in graph 2(REFERENCIA) the plot for the theoretical phase of the output voltage which is similar to a normal bandpass filter.

\begin{figure}[H] \centering
\includegraphics[width=0.7\linewidth]{fase.eps}
\caption{Circuit}
\label{fig:circ}
\end{figure}

\begin{figure}[H] \centering
\includegraphics[width=0.7\linewidth]{teoria.eps}
\caption{Circuit}
\label{fig:circ}
\end{figure}
We also computed the low cutoff frequency and the high cut off frequency, and with those two we compute the central frequency using the following equation:

\begin{equation}
\omega_{L}= \frac{1}{R1*C1}
\end{equation}

\begin{equation}
\omega_{H}= \frac{1}{R2*C2}
\end{equation}

\begin{equation}
\omega_{0}= \sqrt{\omega_{L} * \omega_{H}}
\end{equation}

\begin{table}[H]
\centering
\begin{tabularx}{0.6\textwidth} {
  | >{\raggedright\arraybackslash}X
  | >{\raggedleft\arraybackslash}X | }
 \hline
Lower Cut Off Frequency & 4.545455e+03 rad/s\\ \hline
Higher Cut Off Frequency & 9.090909e+03 rad/s\\ \hline
Central Frequency & 6.428243e+03 rad/s\\ \hline

\end{tabularx}
\caption{Frequencies}
\end{table}

The figure of merit was also computed, following the expression presented in section 1.

\begin{table}[H]
\centering
\begin{tabularx}{0.6\textwidth} {
  | >{\raggedright\arraybackslash}X
  | >{\raggedleft\arraybackslash}X | }
 \hline
Cost & 14057.9\\ \hline
Merit & 0.000269578\\ \hline

\end{tabularx}
\caption{Cost and merit}
\end{table}


\subsection{Theoretical Results}
In the graphic below we present the voltage gain of the circuit, has a function of log10(frequency).
\par

GRAFICO DO VOLTAGE GAIN

\par
In the following tables you can see the values of the input and output impedances; the low cut-off frequency, high cut-off frequency and central frequency; the gain for the central frequency and the merit obtained and the cost of the circuit.

\par
TABELA DE INPUT E OUTPUT IMPEDANCES

\par
TABELA DE CUT OFF FREQ, HOGH CUT OFF FREQ E CENTRAL FREQ

\par
TABELA CENTRAL FREQ E GAIN RESPETIVO

\par
TABELA CUSTO E MERITO



\section{Simulation Analysis}

In this section of the report, we are going to simulate the modeled circuit using Ngspice. Our main objective with this simulation is confirming the validity of our theoretical approach and try to explain any major discrepancies. Therefore, we will be paying close attention to the values of the gain, central frequency and input/output impedances, trying to obtain the best results possible for our application. Finally, we will compute the merit figure, looking forward to having the best possible value.

Firstly, it is important to analyse the frequency domain and determine the output voltage gain. For this, we utilized the .meas function.
The main goal of our work was designing a passband filter, that would cut both low and very high frequencies, obtaining the following results, which are shown in the table below (REFERENCIA).

\begin{table}[H]
\centering
\begin{tabularx}{0.6\textwidth} {
  | >{\raggedright\arraybackslash}X
  | >{\raggedleft\arraybackslash}X | }
 \hline
Central Frequency (Hz)& 1000\\ \hline
Gain (dB)& 39.977\\ \hline
Central frequency deviation&0\\ \hline
Gain deviation&0.263873\\ \hline

\end{tabularx}
\caption{Ngspice results}
\end{table}

\subsection{Impedances}

Now, it's important to take a look at the input and output impedances, as they play a key role in our circuit.

Firstly, looking at our input impedance (REFERENCIA), we wanted to have a very high resistance, as it would imply that the voltage in node 2 is as similiar to the input voltage as possible. This would be very beneficial for the gain, because having a high quocient of $\frac{R_2}{R_3}$, improves significantly the gain value.

On the other hand, for the output impedance, we need to make it as low as possible in order to have the highest possible output voltage. It is clear why if we take a look at the voltage divider formula.

\begin{table}[H]
\centering
\begin{tabularx}{0.6\textwidth} {
  | >{\raggedright\arraybackslash}X
  | >{\raggedleft\arraybackslash}X | }
 \hline
Zin & -1445.55 + 303.662 j\\ \hline

\end{tabularx}
\caption{Input impedance from Ngspice}
\end{table}

\begin{table}[H]
\centering
\begin{tabularx}{0.6\textwidth} {
  | >{\raggedright\arraybackslash}X
  | >{\raggedleft\arraybackslash}X | }
 \hline
Zo & 0.20916 + -14.4657 j\\ \hline

\end{tabularx}
\caption{Output impedance from Ngspice}
\end{table}



Finally, we plotted the merit figure and cost in order to confirm how effective our circuit design was.


\begin{table}[H]
\centering
\begin{tabularx}{0.6\textwidth} {
  | >{\raggedright\arraybackslash}X
  | >{\raggedleft\arraybackslash}X | }
 \hline
Cost & 14057.9\\ \hline
Merit & 0.000269578\\ \hline

\end{tabularx}
\caption{Cost and Merit}
\end{table}
\section{Theoretical Results and Simulation Comparison}

For this section, we are going to perform a general comparison between theoretical results obtained through Octave and simulation results obtained with Ngspice.

Both impedances are similar in theoretical and simulation
analysis, because the OP-AMP is not perfect and ngspice uses a more complex models to analyse it.
Then looking to the gain response, we can see that both of them are similar, but, again, since Ngspice uses more complex models to obtain the gain, we find some differences between them, as we can see in the central frequency and the maximum gain value.

Comparing the phase plots of both analysis, in octave we find two roots and two poles because of the presence of two capacitors, but in Ngspice we have two roots and four poles. This is due to the OP-AMP models used. In octave we considered it an ideal one, but in fact the OP-AMP has capacitors that produce two more poles and, again, Ngspice uses more complex models.

\section{Conclusion}
\label{resan}
In this laboratory assignment we were able to deepen our knowledge regarding BJT's and its applicability in various devices, in our case, an Audio Amplifier. 
As a conclusion we can state that the results of Ngspice and Octave differ. This was expected as we used non-linear components, the transistors, in the circuit and they were replaced in the incremental analysis by two resistances and a dependent current source (linear components).
To conclude, we believe that the main goal of the task proposed was achieved as we reached a high gain and large enough bandwidth to cover the human earing capacity.  Overall this solution was a very successful one.  








\end{document}

